% from ega4-8-fr.tex
\oldpage[IV-3]{32}
de présentation finie sur $S_\alpha$). Soit
$f_\alpha:X_\alpha\to Y_\alpha$ un $S_\alpha$-morphisme. Pour que
$f:X\to Y$ soit un isomorphisme, il faut et il suffit qu'il existe
$\lambda\geq\alpha$ tel que $f_\lambda:X_\lambda\to Y_\lambda$ soit un
isomorphisme.
\end{corollary}

\begin{proof}
La condition est évidemment suffisante. Pour prouver qu'elle est nécessaire,
notons d'abord que la question étant locale sur $S_0$ (puisque $L$ est
filtrant), on peut toujours supposer $S_0$ affine, donc quasi-séparé. Il y a
par hypothèse un $S$-morphisme $g:Y\to X$ tel que $g\circ f=1_X$ et
$f\circ g=1_Y$. Comme $X_\alpha$ est de présentation finie sur $S_\alpha$,
et $Y_\alpha$ quasi-compact et quasi-séparé (puisque $S_\alpha$ est
quasi-compact et quasi-séparé), il existe un $\mu\geq\alpha$ et un
$S_\mu$-morphisme $g_\mu:Y_\mu\to X_\mu$ tels que
$g=g_\mu\times1_S$ en vertu de (8.8.2, (i)). D'autre part, il résulte aussi de
(8.8.2, (i)) et des relations $g\circ f=1_X$ et $f\circ g=1_Y$ qu'il existe
$\nu\geq\mu$ tel que l'on ait $g_\nu\circ f_\nu=1_{X_\nu}$ et
$f_\nu\circ g_\nu=1_{Y_\nu}$, puisque $X_\alpha$ et $Y_\alpha$ sont de type
fini sur $S_\alpha$ et quasi-compacts. Cela signifie que
$f_\nu:X_\nu\to Y_\nu$ est un isomorphisme, d'où le corollaire.
\end{proof}

\begin{corollary}[8.8.2.5]
\label{IV.8.8.2.5-fr}
Supposons $S_0$ quasi-compact et quasi-séparé, $X_\alpha$ et $Y_\alpha$ de
présentation finie sur $S_\alpha$. Pour que $X$ et $Y$ soient
$S$-isomorphes, il faut et il suffit qu'il existe $\lambda\geq\alpha$ tels que
$X_\lambda$ et $Y_\lambda$ soient $S_\lambda$-isomorphes. De plus, pour tout
$S$-isomorphisme $f:X\xrightarrow{\sim}Y$, il existe $\mu\geq\lambda$ et un
$S_\mu$-isomorphisme $f_\mu:X_\mu\to Y_\mu$ tels que
$f=f_\mu\times1_S$.
\end{corollary}

\begin{proof}
La condition est évidemment suffisante~; inversement, s'il existe un
$S$-isomorphisme $f:X\to Y$, il résulte de (8.8.2, (i)) que $f$ est de la
forme $f_\mu\times1_S$ pour un $\mu\geq\alpha$ et un homomorphisme
$f_\mu:X_\mu\to Y_\mu$~; mais comme $f$ est un isomorphisme, il résulte de
(8.8.2.4) qu'il existe $\nu\geq\mu$ tel que
$f_\nu:X_\nu\to Y_\nu$ soit un isomorphisme.
\end{proof}

\begin{proof}[Démonstration de (8.8.2), (ii)]
Considérons encore en premier lieu le cas où $S_0=\operatorname{Spec}(A_0)$ et
$X=\operatorname{Spec}(B)$ sont affines~; alors l'assertion de (ii) équivaut
au lemme (1.8.4.2).

Pour démontrer (ii) dans le cas général, notons que $S$ est quasi-compact et
quasi-séparé~; comme $X$ est de présentation finie sur $S$ et $S$ affine sur
$S_0$, il existe donc un recouvrement fini $(U_i)$ de $S_0$ et, si $(W_i)$ est
le recouvrement ouvert affine de $S$ formé des images réciproques des $U_i$
par le morphisme $S\to S_0$, un recouvrement fini $(X_r)$ de $X$ par des
ouverts affines tel que l'image par $X\to S$ de chaque $X_r$ soit contenue
dans un $W_{i(r)}$~; l'anneau $A(X_r)$ est alors une algèbre de présentation
finie sur l'anneau $A(W_{i(r)})$ (1.4.6). En vertu du lemme (1.8.4.2) et du
fait que $L$ est filtrant, il existe un indice $\lambda\in L$ et, pour chaque
indice $r$, un schéma affine $Z_{r\lambda}$, de présentation finie sur l'image
réciproque $W_{i(r),\lambda}$ de $U_{i(r)}$ dans $S_\lambda$, et un
$S$-isomorphisme
$g_r:Z_{r\lambda}\times_{S_\lambda}S\xrightarrow{\sim}X_r$. Soit $Z_{rs}$
l'image réciproque par $g_r$ du sous-préschéma induit sur l'ouvert
$X_r\cap X_s$ de $X_r$, qui est quasi-compact puisque $X$ est quasi-séparé,
et désignons par $g'_{rs}$ la restriction
$Z_{rs}\xrightarrow{\sim}X_r\cap X_s$ de l'isomorphisme $g_r$. En vertu de
(8.8.2.5), il existe un $\mu\geq\lambda$ et, pour tout couple $(r,s)$, un
ouvert quasi-compact $Z_{rs\mu}$ de
$Z_{r\mu}=v_{\lambda\mu}^{-1}(Z_{r\lambda})$ tels que $Z_{rs}$ soit l'image
réciproque de $Z_{rs\mu}$~; d'ailleurs, comme $S_\mu$ est quasi-séparé, et
$W_{i(r),\mu}$ ouvert quasi-compact dans $S_\mu$, chacun des $Z_{rs\mu}$ est
de présentation finie sur $S_\mu$ (I, 1.6.2). Considérons alors, pour tout
couple $(r,s)$, l'isomorphisme
$h_{sr}=g_{sr}^{'-1}\circ g'_{rs}$ de $Z_{rs}$ sur $Z_{sr}$~; il résulte de
(8.8.2.4) qu'il existe un $\nu\geq\mu$ et, pour tout couple $(r,s)$, un
isomorphisme $h_{sr\nu}:Z_{rs\nu}\to Z_{sr\nu}$ tels que
$h_{sr}=h_{sr\nu}\times1_S$. Enfin, pour tout triplet $(r,s,t)$ désignons par
$h'_{sr}$ la restriction de $h_{sr}$ à
